\section{The 5 Kinds of Business a Founder Might Build}

\subsection{Archetypes of Enterprise}

People speak of starting a business as though it describes 1 economic activity. In practice, it covers organizations with profoundly different objectives:

\begin{enumerate}[label=\textbf{\arabic*.},leftmargin=11pt,itemsep=1.5pt,topsep=1.5pt,parsep=0pt]
  \item \textbf{A self-employment practice} sells the founder's personal expertise or labor. A skilled consultant, independent developer, trainer, designer, or physician may earn well and enjoy autonomy. But if most customers are paying for the founder specifically, the operation remains heavily tied to the founder's availability. This is not a moral failure. It is simply a form of work with limited transferable enterprise value unless systems and relationships develop beyond the individual.
  \item \textbf{A profitable operating company} develops customers, employees, processes, and assets that can produce cash flow without every unit of value being personally delivered by the founder. A manufacturing firm, agency, software company, training business, logistics operator, or specialty retailer can fit this description. Such companies may pursue moderate growth, substantial ownership retention, and durable profit rather than spectacular valuation.
  \item \textbf{A venture-scale startup}, in the sense developed in Graham's \textit{Startup = Growth}, is designed to grow quickly because it serves (or hopes to serve) a very large market through a model capable of rapid expansion.\footnote{Paul Graham, ``Startup = Growth'' (September 2012).} External equity financing may be useful, but it is not the defining characteristic. The defining characteristic is the search for and exploitation of a growth opportunity whose scale could justify extraordinary risk and investment.
  \item \textbf{A capital-intensive enterprise} may require factories, research facilities, land, regulatory approvals, large inventories, or physical infrastructure before it can serve customers at scale. Such businesses can become formidable, but their early constraints are different. Even an excellent idea may need access to specialized financing or staged partnerships before the company can prove commercial viability.
  \item \textbf{An institution} is an organization in which the central asset is not merely a profitable product, but a continuing system of standards, talent formation, credibility, governance, knowledge, and relationships. Universities, laboratories, hospitals, research foundations, and certain enduring companies have institutional characteristics. An institution may have commercial businesses, charitable activities, public responsibilities, or some combination. It may be difficult to measure its success using a single shareholder-value metric.
\end{enumerate}

\subsection{Evolution and Distinct Games}

These categories overlap. A founder-led service can evolve into a profitable operating company. A software startup can become an institution. A research institute can create valuable commercial subsidiaries. A consumer brand can move from local production to mass distribution. The important lesson is that the same founding activity can lead to different kinds of organizations, and the founder should understand which evolution they want.

A profitable small company is not an unsuccessful venture startup. A venture startup that cannot grow rapidly may be an unsuccessful venture investment even while remaining useful to customers. A university with outstanding research and weak finances may be academically important yet financially fragile. A firm generating good shareholder returns through practices that harm its customers or society may be economically successful for its owners without being worthy of admiration.

\textbf{Success must be defined relative to the game being played.}

Research on small-business owners reinforces this distinction. Erik Hurst and Benjamin Wild Pugsley found that many new small-business owners in the United States intended to supply existing services to existing markets and had little intention of growing large. Independence and flexibility were important motivations.\footnote{Erik Hurst and Benjamin Wild Pugsley, ``What Do Small Businesses Do?'' NBER Working Paper 17041 (2011).} These founders are not necessarily trying to create an exponential startup and failing to do so. Many never wanted that game.

If your aim is a financially secure independent life, a high-quality small business may be a better instrument than a venture-backed software company. If your aim is a laboratory capable of financing risky scientific experiments, a narrow professional practice may produce useful seed capital but may not itself finance the eventual ambition. If your aim is to build a durable regional employer, a stable operating firm may be precisely the right result.

\textbf{A person can spend a decade pursuing an objective they never consciously chose. The first defense is to name the game.}
